\setchapterpreamble[u]{\margintoc}
\chapter{Сравнение алгоритмов КС-достижимости}
\label{chpt:CFPQ_Comparison}
\tikzsetfigurename{CFPQ_Comparison_}

В данной главе приводится сравнение алгоритмов КС-достижимости, рассмотренных в предыдущих главах (гл.~\ref{chpt:CFPQ_CYK}--\ref{chpt:GLR}, а также гл.~\ref{chpt:MatrixBasedAlgos}--\ref{chpt:MCFLPQ}).
Обсуждаются рекомендации по выбору алгоритма в зависимости от постановки задачи, а также практические аспекты, такие как возможность параллелизации и производительность на реальных данных.

\section{Рекомендации по выбору алгоритма}

TODO: Критерии выбора: от одного источника или для всех пар, необходимость всех путей или только достижимости, размер грамматики, размер графа и т.д.

\section{Возможность параллелизации}

TODO: Сравнение алгоритмов с точки зрения возможности параллелизации. Матричные методы распараллеливаются легко (библиотеки линейной алгебры, GPGPU). GLL/GLR~--- сложнее.

\section{Сравнение на реальных данных}

TODO: Таблицы и графики сравнения времени работы и потребления памяти на реальных наборах данных (например, графы из статического анализа кода, графовые БД).
